\documentclass{article}
\usepackage[utf8]{inputenc}
\usepackage{amsmath}
\usepackage{geometry}
\geometry{margin=2cm}
\begin{document}

\section*{Questão 143 — ENEM 2024 — Matemática}

Uma casa de shows terá um evento cujo custo total de produção é de R\$ 34\,350,00, sendo que comporta 500 pessoas. O preço do ingresso será de R\$ 130,00 e, normalmente, 60\% das pessoas adquirem meia-entrada, pagando R\$ 65,00 pelo ingresso. Além do faturamento proveniente da venda de ingressos, a casa de shows vende, com 60\% de lucro, bebidas e petiscos ao público no dia do evento.

Após ter vendido todos os 500 ingressos, constatou-se que a quantidade de meias-entradas vendidas superou em 50\% o que estava previsto, impactando o faturamento estimado com a venda de ingressos.

No dia do evento, decidiu-se manter o percentual de 60\% de lucro sobre as bebidas e petiscos, pois todo o público que comprou ingresso compareceu ao show. Com isso, espera-se ter lucro de R\$ 17\,000,00 nesse evento.

\vspace{0.5em}

\textbf{Pergunta:} Para que se alcance o lucro esperado, o gasto médio por pessoa com bebidas e petiscos, em real, deverá ser de?

\section*{Solução Técnica}

A estratégia adotada para resolver o problema foi:

\begin{itemize}
  \item Calcular a quantidade prevista e a quantidade real de meias-entradas vendidas.
  \item Identificar as quantidades de ingressos integrais e meia-entrada efetivamente vendidas.
  \item Calcular a receita real proveniente da venda de ingressos.
  \item Determinar o lucro obtido com a venda de ingressos com base no custo total do evento.
  \item Usar o lucro desejado total para calcular o lucro necessário da venda de bebidas e petiscos.
  \item Aplicar a taxa de lucro de 60\% dessas vendas para calcular o faturamento total com bebidas e petiscos.
  \item Dividir o faturamento total pelas 500 pessoas para obter o gasto médio por pessoa.
\end{itemize}

A resposta final é \textbf{R\$ 52,00} para o gasto médio por pessoa com bebidas e petiscos, que corresponde à alternativa D do gabarito.

\section*{Formulas Utilizadas}

\begin{align*}
&\text{Quantidade prevista de meias} = 60\% \times 500 = 300 \\[6pt]
&\text{Quantidade real de meias} = 300 + 0{,}5 \times 300 = 450 \\[6pt]
&\text{Ingressos integrais} = 500 - 450 = 50 \\[6pt]
&\text{Receita ingressos} = (50 \times 130) + (450 \times 65) = 6.500 + 29.250 = 35.750 \\[6pt]
&\text{Lucro ingressos} = 35.750 - 34.350 = 1.400 \\[6pt]
&\text{Lucro total desejado} = 17.000 \\[6pt]
&\text{Lucro em bebidas e petiscos} = 17.000 - 1.400 = 15.600 \\[6pt]
&\text{Lucro} = 60\% \times \text{Faturamento}, \quad \Rightarrow \quad \text{Faturamento} = \frac{15.600}{0{,}6} = 26.000 \\[6pt]
&\text{Gasto médio por pessoa} = \frac{26.000}{500} = 52
\end{align*}

\section*{Explicação Didática}

\subsection*{Resumo}

Para calcular o gasto médio por pessoa com bebidas e petiscos, considera-se o lucro obtido com ingressos, o custo total do evento e o lucro desejado. O aumento na venda das meias-entradas altera o faturamento dos ingressos, afetando o valor restante para bebidas e petiscos. A partir do lucro esperado e do percentual de lucro na venda de bebidas e petiscos, determina-se o gasto médio necessário por pessoa.

\subsection*{Passos Detalhados}

\begin{enumerate}
  \item Calcule o número previsto de meias-entradas: $60\%$ de 500, que é 300.
  \item Calcule a nova quantidade real de meias-entradas: 50\% a mais, ou seja, $300 + 0{,}5 \times 300 = 450$.
  \item Calcule o número de ingressos integrais: $500 - 450 = 50$.
  \item Calcule a receita da venda de ingressos:
  \[
  (50 \times 130) + (450 \times 65) = 6.500 + 29.250 = 35.750
  \]
  \item Calcule o lucro proveniente dos ingressos:
  \[
  35.750 - 34.350 = 1.400
  \]
  \item Encontre o lucro necessário para bebidas e petiscos, subtraindo o lucro com ingressos do lucro total esperado:
  \[
  17.000 - 1.400 = 15.600
  \]
  \item Considerando 60\% de lucro sobre vendas de bebidas e petiscos, determine o faturamento necessário:
  \[
  \text{Faturamento} = \frac{15.600}{0{,}6} = 26.000
  \]
  \item Calcule o gasto médio por pessoa:
  \[
  \frac{26.000}{500} = 52
  \]
  \item Portanto, o gasto médio necessário é R\$ 52,00.
\end{enumerate}

\subsection*{Erros comuns}

\begin{itemize}
  \item Confundir lucro com faturamento em bebidas e petiscos.
  \item Não considerar o aumento real da quantidade de meia-entrada vendida.
  \item Calcular gasto médio usando apenas o custo, ignorando o lucro.
  \item Ignorar o lucro obtido com ingressos na conta final.
  \item Dividir diretamente o lucro pelo número de pessoas sem cálculos intermediários.
\end{itemize}

\subsection*{Dicas para ensino}

\begin{itemize}
  \item Explique a distinção entre lucro e faturamento.
  \item Use exemplos simplificados para impactar o entendimento dos percentuais.
  \item Reforce a atenção à leitura dos dados do problema.
  \item Ensine a organizar os dados antes de executar cálculos.
  \item Pratique com situações reais similares envolvendo porcentagem e lucro.
\end{itemize}

\section*{Distratores e Justificativas}

\begin{itemize}
  \item \textbf{19,50}: Provável erro por considerar apenas custo ou preço meia-entrada como gasto médio; não considera lucro.
  \item \textbf{28,80}: Erro ao não ajustar o aumento de meia-entrada nem estimar lucro (60\%) corretamente.
  \item \textbf{34,00}: Confusão na média ponderada ou em percentuais aplicados ao custo e lucro.
  \item \textbf{68,70}: Superestimação por confundir lucro e faturamento ou aplicar incorreto o percentual de lucro.
\end{itemize}

\section*{Perfil da Questão}

\begin{itemize}
  \item \textbf{Habilidade principal:} Análise quantitativa e financeira.
  \item \textbf{Habilidades secundárias:} Porcentagem, raciocínio proporcional e cálculo com números reais.
  \item \textbf{Demanda cognitiva:} Média — manipula percentuais financeiros e requer interpretação do problema.
  \item \textbf{Dificuldade:} Média.
  \item \textbf{Observações:} Exige compreensão do conceito de lucro versus faturamento e manipulação de percentuais sobre quantidades e valores financeiros.
\end{itemize}

\section*{Revisão de Consistência}

O conteúdo técnico está correto e alinhado com os dados do enunciado. A explicação didática é clara e organizada, contemplando o raciocínio necessário. As justificativas para as alternativas incorretas esclarecem possíveis erros conceituais. Não foram identificados problemas metodológicos ou conceituais na análise apresentada.

\end{document}